% Supplementary material for the Neurocomputing submission.
\documentclass[preprint,12pt]{elsarticle}
\usepackage{graphicx}
\usepackage{caption}
\usepackage{booktabs}
\usepackage{float}
\usepackage{longtable}
\begin{document}
\renewcommand{\thefigure}{S\arabic{figure}}
\renewcommand{\thetable}{S\arabic{table}}
\section*{Supplementary Material}
The supplement contains the full versions of the primary tables and the sensitivity analyses needed to audit the calibration curve, the task-structured involution control matrix, the zero-shot comparison, the strict run-disjoint LORO analysis, and the cross-backbone results. Supplementary Figs.~S1--S2 report the engineering-cost and training-loss diagnostics cited in Methods; Fig.~S3 visualizes the primary gate-release diagnostic; Fig.~S4 gives the detailed cross-backbone profiles.

\begin{figure}[H]
\centering
\includegraphics[width=\textwidth]{Supplementary_Fig_S2_efficiency_performance_cost.pdf}
\caption{Engineering cost and adaptation-time comparison. The timing analysis uses three subjects per backbone.}
\end{figure}

\begin{figure}[H]
\centering
\includegraphics[width=\textwidth]{Supplementary_Fig_S3_N3_training_validation_loss.pdf}
\caption{Training and inner-validation cross-entropy loss for the tuned ME-pretrained frozen-backbone comparator across the five MI-support fractions.}
\end{figure}

\begin{figure}[H]
\centering
\includegraphics[width=0.94\textwidth]{Fig6_gate_release_diagnostic.pdf}
\caption{Primary gate-release diagnostic. The line shows the median fraction of outer-query trials receiving the private correction, the shaded region shows the interquartile range, and the dashed line shows the prespecified support-derived design target. Release rate is a gate diagnostic rather than a calibrated confidence estimate.}
\end{figure}

\begin{figure}[H]
\centering
\includegraphics[width=0.98\textwidth]{Fig5_cross_backbone_profiles.pdf}
\caption{Detailed cross-backbone transfer profiles. The six panels correspond to the prespecified representation backbones. Solid lines show the EEGMMIDB development protocol and dashed lines show the ordinary five-fold Lajtos 2023 protocol; in each protocol, ReAKT-C is compared with the strongest listed neural comparator for that backbone and support fraction. Bottom ticks show EEGMMIDB fractions and top ticks in the upper row show Lajtos fractions. Missing development cells remain blank and complement the main-text backbone matrices.}
\end{figure}

\clearpage
\setcounter{table}{2}
\begin{table*}[!t]
\caption{Tuned ME-pretrained frozen-backbone classifier for the development cohort. The backbone was frozen and only the EEGNet classifier was optimized. Learning rate was selected from $\{10^{-5},3\times10^{-5},10^{-4}\}$ by inner validation. Values are equal-weight subject means after five outer folds.}
\label{tab:n3_tuned_finetune}
\centering
\scriptsize
\begin{tabular}{c r r r r}
\toprule
MI support & Accuracy (\%) & SD (\%) & Mean selected LR & Mean epochs \\
\midrule
0.2 & 66.65 & 15.73 & 8.43e-05 & 88.3 \\
0.4 & 67.91 & 15.32 & 8.76e-05 & 97.1 \\
0.6 & 68.15 & 16.25 & 7.16e-05 & 96.6 \\
0.8 & 68.56 & 15.56 & 7.43e-05 & 98.1 \\
1.0 & 69.07 & 15.81 & 7.68e-05 & 103.4 \\
\bottomrule
\end{tabular}
\end{table*}
\clearpage
\begin{table*}[!t]
\caption{Primary even-branch component contrast under the formal primary development protocol (EEGNet, 104 participants, MI support fraction 0.20). Five outer folds were averaged within participant before equal-weight participant aggregation.}
\label{tab:component_ablation_original}
\centering
\scriptsize
\resizebox{\textwidth}{!}{%
\begin{tabular}{c r r r r r r r r}
\toprule
MI support & ReAKT-C (\%) & Odd-only (\%) & Gain (pp) & 95\% CI (pp) & Wilcoxon $p$ & FDR $q$ & paired $t$ $p$ & wins/ties/losses\\
\midrule
0.20 & 76.85 & 76.27 & +0.58 & [0.11,1.06] & 0.0284 & 0.0569 & 0.0188 & 37/44/23\\
\bottomrule
\end{tabular}%
}
\end{table*}

\clearpage
\begin{table*}[!t]
\caption{Task-structured involution and target-label controls under the canonical primary protocol (EEGNet, 104 participants, MI support fraction 0.20). The compact main-text table reports accuracy and paired accuracy contrasts; the remaining metrics are given in the supplementary all-metrics table. Accuracy contrasts are ReAKT-C minus the listed variant. Group-only controls are left without imputed paired $p$ values.}
\label{tab:primary_controls}
\centering
\small
\resizebox{\textwidth}{!}{%
\begin{tabular}{l r r c c}
\toprule
Variant & Accuracy (\%) & $\Delta$ (pp) & 95\% CI (pp) & $p$; W/T/L \\
\midrule
ReAKT-C & 76.85 & -- & -- & -- \\
ME-only zero-shot & 77.26 & -0.40 & [-1.42,0.60] & 0.312; 44/17/43 \\
Odd-only & 76.27 & +0.58 & [0.11,1.06] & 0.0284; 37/44/23 \\
Frozen $h$ + RBF-KRR & 72.48 & +4.37 & [2.92,5.78] & $4.17\times10^{-7}$; 70/12/22 \\
Random involution & 73.60 & +3.26 & [1.76,4.74] & $7.53\times10^{-5}$; 63/14/27 \\
Direct $[h,h_{\tau}]$ double-kernel & 72.98 & +3.88 & [2.52,5.18] & $1.61\times10^{-6}$; 68/15/21 \\
ME-only odd + support-even & 76.79 & +0.06 & -- & -- \\
Even-only & 57.62 & +19.23 & -- & -- \\
Odd + even, no gate & 76.35 & +0.50 & -- & -- \\
Local adjacent involution & 73.51 & +3.34 & -- & -- \\
Mirror augmentation & 75.90 & +0.95 & -- & -- \\
\bottomrule
\end{tabular}%
}
\end{table*}

\clearpage
\begin{table*}[!t]
\caption{Secondary metrics for the canonical EEGMMIDB control matrix (EEGNet, MI support fraction 0.20). These values correspond to the compact accuracy table in the main text.}
\label{tab:primary_controls_metrics}
\centering
\scriptsize
\resizebox{\textwidth}{!}{%
\begin{tabular}{l r r r r r}
\toprule
Variant & Accuracy (\%) & Balanced accuracy (\%) & Macro-F1 (\%) & Kappa & ROC-AUC (\%) \\
\midrule
ReAKT-C & 76.85 & 76.72 & 75.21 & 0.529 & 84.15 \\
ME-only zero-shot & 77.26 & 77.29 & 76.12 & 0.539 & 83.33 \\
Odd-only & 76.27 & 76.09 & 74.58 & 0.516 & 83.97 \\
Frozen $h$ + RBF-KRR & 72.48 & 72.23 & 70.01 & 0.440 & 78.81 \\
Random involution & 73.60 & 73.45 & 71.56 & 0.464 & 80.37 \\
Direct $[h,h^R]$ double-kernel & 72.98 & 72.72 & 70.58 & 0.450 & 78.93 \\
ME-only odd + support-even & 76.79 & 76.77 & 75.70 & 0.529 & 82.99 \\
Even-only & 57.62 & 57.18 & 52.72 & 0.142 & 61.57 \\
Odd + even, no gate & 76.35 & 75.95 & 74.50 & 0.516 & 84.25 \\
Local adjacent involution & 73.51 & 73.45 & 71.47 & 0.463 & 80.88 \\
Mirror augmentation & 75.90 & 75.88 & 74.73 & 0.512 & 82.83 \\
\bottomrule
\end{tabular}%
}
\end{table*}

\clearpage
\begin{table*}[!t]
\caption{Canonical EEGMMIDB EEGNet support curve (104 participants, five matched outer folds). ME-only zero-shot is constant because it uses no MI labels. CIs are subject-level bootstrap intervals; $q$ is BH-FDR adjusted over the five paired comparisons, and W/T/L are ReAKT-C wins, ties, and losses. Differences use unrounded means before display rounding.}
\label{tab:zero_shot_curve}
\centering
\scriptsize
\resizebox{\textwidth}{!}{%
\begin{tabular}{c r r r r r c c c c}
\toprule
Support & ReAKT-C (\%) & Odd-only (\%) & ME-only (\%) & ME odd + support-even (\%) & $\Delta$ (pp) & 95\% CI & $p$ & $q$ & W/T/L\\
\midrule
0.20 & 76.85 & 76.27 & 77.26 & 76.79 & -0.40 & [-1.43,+0.63] & 0.312 & 0.312 & 44/17/43\\
0.40 & 79.04 & 78.91 & 77.26 & 76.74 & +1.78 & [+0.96,+2.67] & $8.83\times10^{-4}$ & $1.10\times10^{-3}$ & 52/21/31\\
0.60 & 79.76 & 79.79 & 77.26 & 77.05 & +2.50 & [+1.43,+3.64] & $1.79\times10^{-4}$ & $4.47\times10^{-4}$ & 56/24/24\\
0.80 & 80.26 & 80.36 & 77.26 & 77.04 & +3.00 & [+1.93,+4.13] & $3.62\times10^{-6}$ & $1.81\times10^{-5}$ & 57/20/27\\
1.00 & 79.75 & 79.86 & 77.26 & 77.25 & +2.49 & [+1.33,+3.66] & $3.38\times10^{-4}$ & $5.64\times10^{-4}$ & 58/16/30\\
\bottomrule
\end{tabular}%
}
\end{table*}

\clearpage
\begin{table*}[!t]
\caption{Checkpoint/split audit of the random-involution control at EEGNet and MI support 0.20. The diagnostic exports used distinct audit states and are reported only through their within-export ReAKT-C minus random-involution contrasts. They are an implementation audit, not a seed-controlled estimate of the canonical primary accuracy, and are not pooled with or used to overwrite the canonical result in Table~3.}
\label{tab:random_involution}
\centering
\scriptsize
\resizebox{\textwidth}{!}{%
\begin{tabular}{r r r r}
\toprule
Seed & $\Delta$ vs random (pp) & 95\% CI vs random & $p$ vs random\\
\midrule
20261108 & +2.74 & [1.43,4.05] & 6.57e-05 \\
20261109 & +3.35 & [1.90,4.78] & 1.75e-05 \\
20261110 & +2.89 & [1.50,4.27] & 5.85e-05 \\
20261111 & +4.82 & [3.37,6.25] & 4.10e-08 \\
20261112 & +3.68 & [2.37,5.03] & 2.34e-06 \\
20261113 & +3.46 & [1.90,4.96] & 2.79e-05 \\
20261114 & +4.54 & [3.06,6.01] & 2.22e-07 \\
20261115 & +2.59 & [1.44,3.76] & 2.74e-05 \\
20261116 & +3.04 & [1.61,4.43] & 4.91e-05 \\
20261117 & +2.23 & [0.89,3.56] & 1.53e-03 \\
\bottomrule
\end{tabular}%
}
\end{table*}

\clearpage
\begin{table*}[!t]
\caption{Primary quantile-matched gate release summaries from 520 matched outer folds per support level. The reported fraction is the realized proportion of outer-query trials receiving the private correction; it is compared with the prespecified support-derived release target. The fixed-threshold diagnostic is outside the primary release summary.}
\label{tab:gate_distribution}
\centering
\scriptsize
\begin{tabular}{c r c}
\toprule
MI support & Primary query median [IQR] & Design target $q_{\mathrm{in}}$ \\
\midrule
0.20 & 0.111 [0.111,0.333] & 0.108 \\
0.40 & 0.111 [0,0.222] & 0.113 \\
0.60 & 0.111 [0,0.222] & 0.116 \\
0.80 & 0.111 [0,0.222] & 0.113 \\
1.00 & 0.111 [0,0.222] & 0.111 \\
\bottomrule
\end{tabular}
\end{table*}

\clearpage
\begin{table}[H]
\caption{Paired comparisons of ReAKT-C against ME-only zero-shot transfer. CIs are subject-level bootstrap intervals; $p$ is the raw paired Wilcoxon value and $q$ is BH-FDR adjusted within each task-specific support family. Wins/ties/losses refer to the ReAKT-C minus zero-shot accuracy contrast.}
\label{tab:zero_shot_inference}
\centering\scriptsize
\resizebox{\textwidth}{!}{%
\begin{tabular}{l c r r r r r r c}
\toprule
Task & Support & ReAKT-C (\%) & Zero-shot (\%) & $\Delta$ (pp) & 95\% CI & $p$ & $q$ & W/T/L\\
\midrule
EEGMMIDB-dev & 0.20 & 76.85 & 77.26 & -0.40 & [-1.43,+0.63] & 3.12e-01 & 3.12e-01 & 44/17/43\\
EEGMMIDB-dev & 0.40 & 79.04 & 77.26 & +1.78 & [+0.96,+2.67] & 8.83e-04 & 1.10e-03 & 52/21/31\\
EEGMMIDB-dev & 0.60 & 79.76 & 77.26 & +2.50 & [+1.43,+3.64] & 1.79e-04 & 4.47e-04 & 56/24/24\\
EEGMMIDB-dev & 0.80 & 80.26 & 77.26 & +3.00 & [+1.93,+4.13] & 3.62e-06 & 1.81e-05 & 57/20/27\\
EEGMMIDB-dev & 1.00 & 79.75 & 77.26 & +2.49 & [+1.33,+3.66] & 3.38e-04 & 5.64e-04 & 58/16/30\\
Held-out & 0.20 & 76.35 & 77.26 & -0.91 & [-1.81,-0.00] & 8.05e-02 & 1.01e-01 & 37/17/50\\
Held-out & 0.40 & 77.93 & 77.26 & +0.67 & [-0.38,+1.74] & 2.16e-01 & 2.16e-01 & 50/11/43\\
Held-out & 0.60 & 78.44 & 77.26 & +1.18 & [+0.12,+2.28] & 1.60e-02 & 2.67e-02 & 53/18/33\\
Held-out & 0.80 & 79.88 & 77.26 & +2.63 & [+1.44,+3.85] & 9.78e-05 & 4.89e-04 & 62/14/28\\
Held-out & 1.00 & 79.84 & 77.26 & +2.59 & [+1.35,+3.86] & 4.34e-04 & 1.09e-03 & 63/8/33\\
Lajtos-5F & 0.25 & 76.02 & 74.84 & +1.18 & [-0.43,+2.74] & 1.15e-01 & 0.157 & 18/1/12\\
Lajtos-5F & 0.50 & 75.97 & 74.84 & +1.13 & [-0.16,+2.47] & 1.18e-01 & 0.157 & 19/0/12\\
Lajtos-5F & 0.75 & 75.54 & 74.84 & +0.70 & [-0.97,+2.26] & 4.59e-01 & 0.459 & 17/3/11\\
Lajtos-5F & 1.00 & 76.40 & 74.84 & +1.56 & [-0.43,+3.39] & 5.45e-02 & 0.157 & 20/0/11\\
\bottomrule
\end{tabular}%
}
\end{table}

\clearpage
\begin{table*}[!t]
\caption{Strict run-disjoint Lajtos 2023 leave-one-run-out (LORO) analysis. All 31 participants contributed three held-out MI runs. Each query run contained 20 trials (10 per class); the remaining two runs supplied 5, 10, 15, or 20 support trials per class at fractions 0.25, 0.50, 0.75, and 1.00. CIs are participant-level percentile-bootstrap 95\% CIs for ReAKT-C minus ME-only zero-shot.}
\label{tab:lajtos_loro}
\centering\scriptsize
\resizebox{\textwidth}{!}{%
\begin{tabular}{c r r r r c c c}
\toprule
Support & ReAKT-C Acc. (\%) & ME-only Acc. (\%) & $\Delta$ (pp) & 95\% CI (pp) & Wilcoxon $p$ & W/T/L & $n$ \\
\midrule
0.25 & 75.16 & 74.84 & +0.32 & [-1.13,1.77] & 0.572 & 13/8/10 & 31 \\
0.50 & 75.70 & 74.84 & +0.86 & [-0.54,2.26] & 0.225 & 17/3/11 & 31 \\
0.75 & 75.86 & 74.84 & +1.02 & [-0.81,2.80] & 0.241 & 16/3/12 & 31 \\
1.00 & 76.29 & 74.84 & +1.45 & [-0.27,3.17] & 0.121 & 16/6/9 & 31 \\
\bottomrule
\end{tabular}%
}
\end{table*}

\clearpage
\begin{table}[H]
\caption{Friedman omnibus tests for the development control family. The family contains nine prespecified methods at each support fraction: complete ReAKT-C, ME-only zero-shot, ME-only odd plus support-even, odd-only, even-only, direct-sum no-gate, direct same-normalization, local-adjacent involution, and mirror augmentation. $q$ is BH-FDR adjusted across the five support-specific omnibus tests.}
\label{tab:friedman_fdr}
\centering\scriptsize
\begin{tabular}{c r r r}
\toprule
Support & Methods ($m$) & Friedman $p$ & BH $q$\\
\midrule
0.20 & 9 & 5.30e-47 & 6.63e-47\\
0.40 & 9 & 5.82e-44 & 5.82e-44\\
0.60 & 9 & 1.08e-50 & 2.70e-50\\
0.80 & 9 & 5.90e-52 & 2.95e-51\\
1.00 & 9 & 2.73e-50 & 4.56e-50\\
\bottomrule
\end{tabular}
\end{table}

\clearpage
\begin{table*}[!t]
\caption{Separate mechanism-control inference from a diagnostic export (Supplementary Table~S13). Only within-export paired contrasts are reported; these diagnostics do not replace the canonical primary component estimate. The controls quantify direct same-normalization, ungated, local-adjacent, mirror-augmentation, and source-only-odd plus support-even variants. CIs are subject-level bootstrap intervals; $p$ is two-sided paired Wilcoxon; W/T/L are complete wins, ties, and losses.}
\label{tab:mechanism_diagnostic}
\centering
\scriptsize
\resizebox{\textwidth}{!}{%
\begin{tabular}{l r r r c c}
\toprule
Variant & Complete$-$variant (pp) & 95\% CI (pp) & $p$ & W/T/L & Export \\
\midrule
Direct double-kernel, same normalization & +3.41 & [2.08,4.73] & $6.36\times10^{-6}$ & 67/13/24 & D4 diagnostic \\
Odd + even, no gate & +0.80 & [-0.03,1.63] & 0.154 & 44/27/33 & D4 diagnostic \\
Local adjacent involution & +3.64 & [2.26,5.02] & $5.05\times10^{-6}$ & 68/14/22 & D4 diagnostic \\
Mirror augmentation & +1.25 & [-0.04,2.54] & 0.075 & 52/15/37 & D4 diagnostic \\
ME-only odd + support-even & +0.36 & [-0.65,1.38] & 0.516 & 47/20/37 & D4 diagnostic \\
\bottomrule
\end{tabular}%
}
\end{table*}

\clearpage
\begin{table*}[!t]
\caption{Mapping of the manuscript's primary and diagnostic exports. The canonical primary export supplies the headline EEGMMIDB component and support-curve estimates, and the strict LORO export supplies the full-cohort external run-disjoint estimates. Diagnostic exports are retained to expose implementation differences and are not pooled with or used to overwrite canonical values.}
\label{tab:export_map}
\centering\tiny
\setlength{\tabcolsep}{2pt}
\resizebox{\textwidth}{!}{%
\begin{tabular}{p{0.19\textwidth} p{0.19\textwidth} p{0.18\textwidth} p{0.16\textwidth} p{0.22\textwidth}}
\toprule
Export & Run / state & Gate / protocol & Complete and comparator & Role in manuscript \\ 
\midrule
Canonical primary & primary protocol (20261108) & Quantile-matched primary gate; five matched EEGMMIDB folds & 76.85 / 76.27 at support 0.20 & Headline component, support curve, primary controls \\ 
Checkpoint/split audit & retained audit exports & Distinct audit checkpoint/split states and runner metadata & Diagnostic contrasts only; no primary accuracy is reported & Audit only; does not update the canonical result \\ 
Legacy fixed-threshold diagnostic & Legacy fixed-threshold gate export & Fixed-threshold release-rate diagnostic; its accuracy contrasts are outside the primary tables & Not used as the primary complete estimate & Release-rate diagnostic outside the primary results \\
LOO-gate mechanism diagnostic & mechanism-control diagnostic export & Support-LOO-score gate; separate mechanism-control execution & Within-export paired contrasts only & Mechanism-control inference only (S13); not a primary performance estimate \\
Strict Lajtos LORO & 20261111 full-cohort rerun; 31 participants & Three run-disjoint folds; 20 query trials/run; support 5/10/15/20 per class & 75.16 / 75.70 / 75.86 / 76.29\newline vs 74.84 zero-shot & Canonical external run-disjoint analysis (S11) \\ 
\bottomrule
\end{tabular}%
}
\end{table*}





\clearpage
\begin{table*}[!t]
\caption{Development accuracy by representation backbone. Values are subject-level mean accuracies in percent; rows are grouped by backbone and MI support fraction. The symbol ``--'' denotes a cell not run under the prespecified allocation. ReAKT-C was numerically highest in all 20 non-missing development cells.}
\label{tab:development_backbone}
\centering
\scriptsize
\resizebox{\textwidth}{!}{%
\begin{tabular}{l c r r r r r r}
\toprule
Backbone & MI support & ME-Finetune & ME-Joint & Deep CORAL & DANN & CDAN & ReAKT-C \\
\midrule
EEGNet & 0.20 & 72.63 & 70.82 & 72.83 & 73.54 & 73.49 & 76.85 \\
EEGNet & 0.40 & 73.93 & 72.66 & 73.85 & 74.66 & 74.62 & 79.04 \\
EEGNet & 0.60 & 75.33 & 74.06 & 75.54 & 75.84 & 75.85 & 79.76 \\
EEGNet & 0.80 & 75.01 & 75.20 & 75.29 & 76.02 & 75.97 & 80.26 \\
EEGNet & 1.00 & 75.62 & 76.15 & 75.41 & 75.83 & 75.81 & 79.75 \\
\addlinespace
EEG-Conformer & 0.20 & 67.94 & 64.86 & 70.08 & 70.10 & 69.74 & 74.22 \\
EEG-Conformer & 0.40 & 70.14 & 66.30 & 71.60 & 71.96 & 72.06 & 75.63 \\
EEG-Conformer & 0.60 & -- & -- & -- & -- & -- & -- \\
EEG-Conformer & 0.80 & -- & -- & -- & -- & -- & -- \\
EEG-Conformer & 1.00 & 73.42 & 72.80 & 73.52 & 74.46 & 74.24 & 76.95 \\
\addlinespace
EEGNeX & 0.20 & 71.99 & 71.50 & 72.42 & 72.61 & 72.31 & 76.39 \\
EEGNeX & 0.40 & 73.02 & 73.13 & 74.12 & 73.49 & 73.45 & 77.51 \\
EEGNeX & 0.60 & -- & -- & -- & -- & -- & -- \\
EEGNeX & 0.80 & -- & -- & -- & -- & -- & -- \\
EEGNeX & 1.00 & 75.56 & 74.94 & 76.66 & 76.36 & 76.27 & 78.71 \\
\addlinespace
EEG-TCNet & 0.20 & 67.75 & 61.70 & 67.86 & 67.83 & 67.93 & 74.80 \\
EEG-TCNet & 0.40 & 68.21 & 65.93 & 69.26 & 69.18 & 69.11 & 75.45 \\
EEG-TCNet & 0.60 & -- & -- & -- & -- & -- & -- \\
EEG-TCNet & 0.80 & -- & -- & -- & -- & -- & -- \\
EEG-TCNet & 1.00 & 71.55 & 70.67 & 72.13 & 71.76 & 71.85 & 77.53 \\
\addlinespace
ATCNet & 0.20 & 71.30 & 67.02 & 72.01 & 71.70 & 71.70 & 75.80 \\
ATCNet & 0.40 & 72.00 & 68.96 & 73.61 & 72.70 & 72.75 & 77.83 \\
ATCNet & 0.60 & -- & -- & -- & -- & -- & -- \\
ATCNet & 0.80 & -- & -- & -- & -- & -- & -- \\
ATCNet & 1.00 & 75.38 & 75.22 & 75.86 & 75.76 & 75.98 & 78.52 \\
\addlinespace
MSVTNet & 0.20 & 63.89 & 61.41 & 65.00 & 65.03 & 65.01 & 66.89 \\
MSVTNet & 0.40 & 64.77 & 62.26 & 65.95 & 65.47 & 65.27 & 67.63 \\
MSVTNet & 0.60 & -- & -- & -- & -- & -- & -- \\
MSVTNet & 0.80 & -- & -- & -- & -- & -- & -- \\
MSVTNet & 1.00 & 66.51 & 65.04 & 67.04 & 67.29 & 67.40 & 67.62 \\
\bottomrule
\end{tabular}
}
\end{table*}


\clearpage
\begin{table*}[!t]
\caption{External Lajtos 2023 accuracy by representation backbone under the ordinary five-fold protocol. Values are subject-level mean accuracies in percent. The displayed columns are the prespecified neural comparators and ReAKT-C; ``--'' denotes cells not run under the allocation.}
\label{tab:external_backbone}
\centering
\scriptsize
\resizebox{\textwidth}{!}{%
\begin{tabular}{l c r r r r r r}
\toprule
Backbone & MI support & ME-Finetune & ME-Joint & Deep CORAL & DANN & CDAN & ReAKT-C\\
\midrule
EEGNet & 0.25 & 70.54 & 62.15 & 70.32 & 70.05 & 70.16 & 76.02\\
EEGNet & 0.50 & 69.52 & 68.44 & 69.84 & 70.81 & 71.13 & 75.97\\
EEGNet & 0.75 & 71.08 & 71.34 & 71.18 & 71.40 & 71.72 & 75.54\\
EEGNet & 1.00 & 71.56 & 75.81 & 72.37 & 71.94 & 71.99 & 76.40\\
\addlinespace
EEG-Conformer & 0.25 & 66.29 & 58.28 & 67.15 & 67.04 & 67.37 & 73.92\\
EEG-Conformer & 0.50 & 68.33 & 65.59 & 71.18 & 68.55 & 68.28 & 73.17\\
EEG-Conformer & 0.75 & 68.87 & 69.46 & 72.10 & 70.48 & 70.16 & 72.90\\
EEG-Conformer & 1.00 & 72.47 & 72.04 & 74.25 & 73.01 & 73.49 & 74.14\\
\addlinespace
EEGNeX & 0.25 & 70.75 & 67.31 & 72.42 & 69.95 & 70.75 & 75.43\\
EEGNeX & 0.50 & 72.15 & 74.09 & 72.80 & 71.34 & 70.91 & 76.61\\
EEGNeX & 0.75 & 72.85 & 77.20 & 74.73 & 72.58 & 73.71 & 77.10\\
EEGNeX & 1.00 & 74.25 & 80.38 & 76.34 & 75.75 & 74.52 & 77.58\\
\addlinespace
EEG-TCNet & 0.25 & 59.03 & 53.55 & 59.68 & 59.57 & 59.41 & 68.39\\
EEG-TCNet & 0.50 & 61.24 & 59.73 & 61.67 & 62.15 & 62.04 & 70.91\\
EEG-TCNet & 0.75 & 62.42 & 64.57 & 64.14 & 65.75 & 65.97 & 72.37\\
EEG-TCNet & 1.00 & 64.35 & 66.13 & 67.37 & 68.01 & 67.74 & 72.90\\
\addlinespace
ATCNet & 0.25 & 66.18 & 59.41 & 66.24 & 66.34 & 66.02 & 72.90\\
ATCNet & 0.50 & 66.67 & 64.14 & 68.55 & 68.98 & 69.09 & 74.09\\
ATCNet & 0.75 & 68.98 & 69.95 & 70.48 & 71.99 & 71.83 & 74.19\\
ATCNet & 1.00 & 70.86 & 72.04 & 71.56 & 72.96 & 72.74 & 74.57\\
\addlinespace
MSVTNet & 0.25 & 60.27 & 54.14 & 61.13 & 59.62 & 59.09 & 65.11\\
MSVTNet & 0.50 & 63.23 & 61.61 & 62.63 & 61.88 & 62.58 & 66.61\\
MSVTNet & 0.75 & 66.88 & 63.92 & 68.39 & 67.26 & 67.37 & 66.83\\
MSVTNet & 1.00 & 67.63 & 69.52 & 68.71 & 70.32 & 70.65 & 68.92\\
\bottomrule
\end{tabular}%
}
\end{table*}



\clearpage
\begin{table*}[!t]
\caption{FBCSP-LDA comparison for the rerun protocols. FBCSP-LDA was fitted only on the labelled MI support trials. Differences are ReAKT-C minus FBCSP-LDA; Lajtos $q$ values are BH-FDR adjusted across its four support comparisons. EEGMMIDB entries report subject-mean 95\% normal intervals.}
\label{tab:fbcsp_inference}
\centering\scriptsize
\resizebox{\textwidth}{!}{%
\begin{tabular}{l c r r r r l}
\toprule
Dataset & Support & $n$ & ReAKT-C (\%) & FBCSP-LDA (\%) & $\Delta$ (pp) & 95\% CI / Wilcoxon $p$ / $q$ \\
\midrule
EEGMMIDB & 0.20 & 104 & 76.85 & 50.97 & +25.88 & [50.00,51.94] / -- / -- \\
EEGMMIDB & 0.40 & 104 & 79.04 & 52.50 & +26.54 & [51.21,53.79] / -- / -- \\
EEGMMIDB & 0.60 & 104 & 79.76 & 53.18 & +26.58 & [51.47,54.90] / -- / -- \\
EEGMMIDB & 0.80 & 104 & 80.26 & 56.09 & +24.17 & [54.21,57.96] / -- / -- \\
EEGMMIDB & 1.00 & 104 & 79.75 & 53.32 & +26.43 & [51.25,55.39] / -- / -- \\
\addlinespace
Lajtos 2023 & 0.25 & 31 & 76.02 & 59.68 & +16.34 & 1.67e-4 / 6.68e-4 \\
Lajtos 2023 & 0.50 & 31 & 75.97 & 64.25 & +11.72 & 3.72e-3 / 7.45e-3 \\
Lajtos 2023 & 0.75 & 31 & 75.54 & 68.17 & +7.37 & 4.35e-2 / 5.80e-2 \\
Lajtos 2023 & 1.00 & 31 & 76.40 & 72.15 & +4.25 & 2.36e-1 / 2.36e-1 \\
\bottomrule
\end{tabular}
}
\end{table*}

\end{document}


